\par\noindent\textbf{(1)} 线化电子方程m\_e v̇=−e(E+v×B0)，取圆偏振v\_±=v\_x±iv\_y、E\_±=E\_x±iE\_y，得到不同有效分母(ω∓Ω\_e)，其中Ω\_e=eB0/m\_e按正大小定义。由J=−n\_e e v和ε\_eff=1+iσ/(ε0ω)，两旋向折射率满足n\_±²=1−ω\_p²/[ω(ω∓Ω\_e)]，ω\_p²=n\_e e²/(ε0m\_e)。
\par\noindent\textbf{(2)} 波数k\_±=ω n\_±/c；传播要求n\_±²>0，否则该偏振为倏逝/截止支。ω≈|Ω\_e|附近一支出现回旋共振分母，理想无碰撞表达式形式发散，实际需考虑阻尼、温度及吸收，不能把其解释为无限可用旋光。
\par\noindent\textbf{(3)} 当ω≫ω\_p,|Ω\_e|，展开n\_±≈1−ω\_p²/(2ω²)(1±Ω\_e/ω)，故Δk=k\_+−k\_-≈−ω\_p²Ω\_e/(cω²)（具体正负依偏振手性约定）。线偏振是两圆偏振等幅叠加，传播距离L后旋转角θ\_F=(k\_+−k\_-)L/2，量级|θ\_F|≈ω\_p²|Ω\_e|L/(2cω²)。
\par\noindent\textbf{(4)} 反转B0令Ω\_e带符号翻转，旋光方向翻转；n\_e→0有ω\_p→0、θ\_F→0，回到真空。高温等离子体有动理学回旋和Landau效应，若ω接近共振或波长接近Debye尺度，冷无碰撞模型不再定量可靠。
